\documentclass[11pt]{article}
\usepackage[margin=1in]{geometry}
\usepackage{amsmath,amssymb,amsthm}
\usepackage{graphicx}
\usepackage[hidelinks]{hyperref}
\setlength{\emergencystretch}{2em}

\newtheorem{theorem}{Theorem}
\newtheorem{lemma}{Lemma}
\newtheorem{corollary}{Corollary}
\newtheorem{definition}{Definition}

\newcommand{\Lcal}{\mathcal L}
\newcommand{\Kcal}{\mathcal K}
\newcommand{\B}{\mathcal B}
\newcommand{\K}{\mathcal K}
\newcommand{\N}{\mathbb N}
\newcommand{\C}{\mathbb C}

\title{A Hilbert-Module Block Criterion for Coronas Failing BCP}
\author{Run \texttt{fa\_banach\_001}, lane 19}
\date{June 17, 2026}

\begin{document}
\maketitle

\section*{Source Problem}

The source paper is Jinghao Huang, Karimbergenov Kudaybergenov, and Rui Liu,
\emph{The ball-covering property of von Neumann algebras and noncommutative
symmetric spaces}, arXiv:2311.14257.  After proving that a commutative unital
\(C^*\)-algebra \(C(K)\) has BCP/UBCP exactly when \(K\) has a countable
\(\pi\)-basis, the authors ask for a noncommutative version:

\begin{center}
\includegraphics[width=0.92\textwidth]{figures/open_problem_crop.png}
\end{center}

\section*{Definitions}

Let \(E\) be a Hilbert \(C^*\)-module.  Write \(\Lcal(E)\) for its adjointable
operators, \(\Kcal(E)\) for its compact operators, and
\[
        Q:\Lcal(E)\longrightarrow \Lcal(E)/\Kcal(E)
\]
for the quotient map.

For a finite set \(F\subseteq\N\), write
\[
        e_F=\sum_{n\in F} e_n.
\]

\begin{definition}
\label{def:block-system}
A sequence \((e_n)\) of nonzero pairwise orthogonal projections in
\(\Kcal(E)\) is called a BCP block system for the corona \(\Lcal(E)/\Kcal(E)\)
if the following hold.

\begin{enumerate}
\item \emph{Compact-vanishing.}  For every \(C\in\Kcal(E)\) and every
\(\varepsilon>0\), there is a finite set \(F_0\subseteq\N\) such that
\[
        \|C e_F\|<\varepsilon
\]
whenever \(F\subseteq\N\setminus F_0\) is finite.

\item \emph{Finite block equivalence.}  Whenever finite sets \(F,G\subseteq\N\)
have the same cardinality, there is a partial isometry \(v\in\Lcal(E)\) with
\[
        v^*v=e_F,\qquad vv^*=e_G.
\]

\item \emph{Strict gluing.}  If \((F_k)\) and \((G_k)\) are pairwise disjoint
finite subsets of \(\N\), \(|F_k|=|G_k|\), and
\[
        v_k^*v_k=e_{F_k},\qquad v_kv_k^*=e_{G_k},
\]
then there is \(V\in\Lcal(E)\), \(\|V\|\le1\), such that
\[
        V e_{F_k}=v_k,\qquad k\ge1.
\]

\item \emph{Tail norm recovery.}  For every \(T\in\Lcal(E)\), every finite
set \(F_0\subseteq\N\), and every \(\varepsilon>0\), there is a nonempty finite
set \(F\subseteq\N\setminus F_0\) such that
\[
        \|T e_F\|>\|Q(T)\|-\varepsilon.
\]
\end{enumerate}
\end{definition}

\section*{Main Result}

\begin{theorem}
\label{thm:block-criterion}
If the corona \(\Lcal(E)/\Kcal(E)\) admits a BCP block system, then
\(\Lcal(E)/\Kcal(E)\) fails BCP.
\end{theorem}

\begin{corollary}
\label{cor:unital-recovered}
Let \(H\) be an infinite-dimensional separable Hilbert space and let \(D\) be a
nonzero unital \(C^*\)-algebra.  Then
\[
        (\B(H)\otimes D)/(\K(H)\otimes D)
\]
fails BCP.
\end{corollary}

\section*{Idea of the Proof}

The block system is designed to isolate exactly the mechanism used in the
stabilized Calkin proof.  Tail norm recovery gives, for each proposed center,
a late compact block on which that center nearly realizes its quotient norm.
Compact-vanishing and finite block equivalence let us move a copy of that
block to a fresh range where the finite-column compact operator \(T e_F\) is
almost invisible.  Strict gluing assembles these local partial isometries into
one adjointable operator.  Its corona class has norm one and remains outside
every ball \(B(\xi,\|\xi\|)\) from the proposed countable cover.

\section*{Proof}

We use the standard equivalent form of BCP: if a normed space \(Y\) has BCP,
then \(S_Y\) is covered by countably many balls \(B(y_j,\|y_j\|)\).

\begin{lemma}[fresh range escape]
\label{lem:fresh-range}
Let \((e_n)\) be a BCP block system.  Let \(F\subseteq\N\) be finite and
nonempty, let \(X\in \Kcal(E)e_F\), let \(F_0\subseteq\N\) be finite, and let
\(\varepsilon>0\).  Then there are a finite set
\[
        G\subseteq\N\setminus F_0,\qquad |G|=|F|,
\]
and a partial isometry \(v\in\Lcal(E)\) such that
\[
        v^*v=e_F,\qquad vv^*=e_G,
\]
and
\[
        \|v-X\|\ge \|X\|-\varepsilon.
\]
\end{lemma}

\begin{proof}
By compact-vanishing applied to \(X^*\), enlarge \(F_0\) if needed so that
\[
        \|e_G X\|=\|X^*e_G\|<\varepsilon
\]
for every finite \(G\subseteq\N\setminus F_0\).  Choose such a \(G\) with
\(|G|=|F|\), and choose \(v\) by finite block equivalence.  Since
\((1-e_G)v=0\),
\[
        \|v-X\|
        \ge \|(1-e_G)(v-X)\|
        =
        \|(1-e_G)X\|
        \ge \|X\|-\|e_GX\|
        >
        \|X\|-\varepsilon.
\]
\end{proof}

\begin{proof}[Proof of Theorem~\ref{thm:block-criterion}]
Let
\[
        \xi_i=Q(T_i),\qquad i\ge1,
\]
be an arbitrary countable family of centers in \(\Lcal(E)/\Kcal(E)\), with
\(T_i\in\Lcal(E)\).  Put \(r_i=\|\xi_i\|\).  Choose a map
\(\pi:\N\to\N\) such that every value occurs infinitely often.

We recursively choose finite domain blocks \(F_k\), finite range blocks
\(G_k\), and partial isometries \(v_k\).  Suppose the previous blocks have
been chosen.  Let \(H_{k-1}\) be the finite union of all previous \(F_j\)'s and
\(G_j\)'s.  Put \(i=\pi(k)\).  By tail norm recovery, choose a nonempty finite
set
\[
        F_k\subseteq\N\setminus H_{k-1}
\]
such that
\[
        \|T_i e_{F_k}\|>r_i-\frac1k.
\]
Apply Lemma~\ref{lem:fresh-range} to
\[
        F=F_k,\qquad X=T_i e_{F_k},
\]
and to a finite forbidden set containing \(H_{k-1}\cup F_k\).  We obtain
\[
        G_k\subseteq \N\setminus(H_{k-1}\cup F_k)
\]
with \(|G_k|=|F_k|\), and a partial isometry \(v_k\) satisfying
\[
        v_k^*v_k=e_{F_k},\qquad v_kv_k^*=e_{G_k},
\]
and
\[
        \|v_k-T_i e_{F_k}\|>r_i-\frac2k.
\]

By strict gluing, there is \(V\in\Lcal(E)\), \(\|V\|\le1\), such that
\[
        V e_{F_k}=v_k,\qquad k\ge1.
\]
We claim that \(\|Q(V)\|=1\).  The upper bound is clear.  For every
\(C\in\Kcal(E)\), compact-vanishing gives
\[
        \|C e_{F_k}\|\longrightarrow0,
\]
and hence
\[
        \|V-C\|
        \ge
        \|(V-C)e_{F_k}\|
        \ge
        \|v_k\|-\|C e_{F_k}\|
        =
        1-\|C e_{F_k}\|.
\]
Taking \(k\to\infty\) and then the infimum over \(C\in\Kcal(E)\) gives
\(\|Q(V)\|\ge1\).

Fix \(i\).  If \(r_i=0\), then
\[
        \|Q(V)-\xi_i\|\ge0=r_i.
\]
If \(r_i>0\), let \(k\to\infty\) along the infinite set
\(\{k:\pi(k)=i\}\).  For any \(C\in\Kcal(E)\),
\[
        \|V-T_i-C\|
        \ge
        \|(V-T_i-C)e_{F_k}\|
        \ge
        \|v_k-T_i e_{F_k}\|-\|C e_{F_k}\|
        >
        r_i-\frac2k-\|C e_{F_k}\|.
\]
Letting \(k\to\infty\) along this subsequence gives
\[
        \|V-T_i-C\|\ge r_i.
\]
Taking the infimum over \(C\in\Kcal(E)\), we get
\[
        \|Q(V)-\xi_i\|\ge r_i=\|\xi_i\|.
\]
Thus \(Q(V)\in S_{\Lcal(E)/\Kcal(E)}\) is outside every ball
\[
        B(\xi_i,\|\xi_i\|).
\]
The countable family of centers was arbitrary, so the corona fails BCP.
\end{proof}

\begin{proof}[Proof of Corollary~\ref{cor:unital-recovered}]
Let \(E=H\otimes D\) be the standard Hilbert \(D\)-module.  Since \(D\) is
unital,
\[
        \Lcal(E)\cong \B(H)\otimes D,\qquad
        \Kcal(E)\cong \K(H)\otimes D.
\]
Fix an orthonormal basis \((h_n)\) of \(H\), and let \(e_n\) be the projection
onto the \(D\)-summand \(h_n\otimes D\).  Then \(e_n\in\Kcal(E)\), the
compact-vanishing property follows from norm approximation by finite matrices
in \(\K(H)\otimes D\), and finite block equivalence is implemented by the
usual matrix-unit partial isometries tensored with \(1_D\).  Strict gluing is
the strong operator sum of those partial isometries, tensored with \(1_D\).

It remains only to recall tail norm recovery.  If \(f\) is a finite sum of
coordinate projections, \(T\in \B(H)\otimes D\), and \(r=\|Q(T)\|\), then
\[
        r\le \|T(1-f)\|.
\]
Represent \(D\) faithfully on a Hilbert space \(K_D\).  Choose a unit vector
\(\xi\in(1-f)H\otimes K_D\) with \(\|T(1-f)\xi\|>r-\varepsilon/2\), and
approximate \(\xi\) by a vector with finite \(H\)-support, close enough to
preserve the weaker lower bound \(r-\varepsilon\).  The corresponding finite
set \(F\) of coordinates is disjoint from the support of \(f\), and
\[
        \|T e_F\|>r-\varepsilon.
\]
Thus \((e_n)\) is a BCP block system, so Theorem~\ref{thm:block-criterion}
applies.
\end{proof}

\section*{Scope and Nonunital Obstruction}

This packet is not merely a restatement of the unital stabilized-corona
theorem.  It isolates the exact projectional block mechanism behind that
proof.  The same argument would apply to any Hilbert-module corona once a BCP
block system is verified.

For a nonunital coefficient algebra \(D\), the standard coordinate projections
on the Hilbert \(D\)-module \(\ell_2(D)\) are generally multiplier projections,
not compact module projections.  Thus compact block projections and tail norm
recovery are not automatic.  This is the current obstruction to promoting the
unital stabilized-corona theorem to arbitrary nonunital stable multiplier
coronas.

\section*{Verification and Search Bounds}

The run indexes were searched for arXiv:2311.14257 together with Hilbert
module, block frame, tail recovery, orthogonal compact projections, multiplier
corona, and BCP.  No existing packet for this Hilbert-module block criterion
was found.  The proof is self-contained given the block-system axioms and uses
only standard Hilbert-module notation.

\section*{Human Review Focus}

Review Definition~\ref{def:block-system}, especially whether strict gluing and
tail norm recovery are stated at the right level of generality.  The proof of
Theorem~\ref{thm:block-criterion} is then a direct diagonal escape argument.

\begin{thebibliography}{9}

\bibitem{HKL}
Jinghao Huang, Karimbergenov Kudaybergenov, and Rui Liu,
\emph{The ball-covering property of von Neumann algebras and noncommutative
symmetric spaces},
arXiv:2311.14257, 2023.

\end{thebibliography}

\end{document}
